% SHARD-ID: AQM-93-FINITE-RING-HAAH-LAYERS
% SHARD-TITLE: Finite-Ring Haah Layers: Residues, Thickenings, and Blockers
% SHARD-SUMMARY: Defines a layer stack from bare finite spectra to Haah-like Weyl data without promoting arbitrary finite rings to codes.
% SHARD-SUMMARY: Separates residue-field Weyl labels from nilpotent-sensitive Frobenius or generating-character layers.
% SHARD-SUMMARY: Uses finite-field, product, mixed-residue, dual-number, and cyclic quotient examples to mark available and blocked structures.
% SHARD-KEYWORDS: finite rings, Haah, Spec R, residue fields, Frobenius rings, Weyl-Heisenberg, Laurent modules, blockers

\section{Finite-Ring Haah Layers: Residues, Thickenings, and Blockers}
\label{sec:finite-ring-haah-layers}

\subsection{Status and Source Anchors}

\textbf{Status: local synthesis and definitions.}  This shard adds no source
manifest, convention entry, script, run bundle, generated data, populated
finite-ring database row, or all-finite-ring representation theorem.  It
uses the layer vocabulary fixed in convention \texttt{(ba)} and the
finite-ring QSA boundaries in AQM-80--AQM-83.  The Haah-side anchors are
AQM-91, convention \texttt{(az)}, convention \texttt{(ba)}, and
\path{references/lattice_stabilizer_codes/SOURCES.md}, lines 17--43 and
87--99, which register the Laurent-module Pauli, stabilizer, excitation-map,
toric-code matrix, and two-dimensional classification sources.  The
finite-Weyl anchors are AQM-71, AQM-78, AQM-31, AQM-36, AQM-37, and
\path{references/heisenberg_weil/SOURCES.md}, especially the finite
Frobenius-ring and generating-character locators for
Gluesing-Luerssen--Pllaha and Sidana--Kashyap.

The scheme-theoretic anchor is
\path{references/algebraic_geometry/SOURCES.md}; its
\path{references/algebraic_geometry/stacks_project_algebra.tex} entry records
prime/maximal ideals, residue fields, spectra, Artinian finite-spectrum
facts, and finite-discrete spectrum equivalences.  The only all-finite-ring
step used below is local: a finite ring has finitely many ideals, so
descending ideal chains stabilize; with the cited Artinian-spectrum source,
its prime spectrum is finite and all primes are maximal.  If \(x\in\Spec R\),
then \(R/x\) is a finite domain, hence a field because multiplication by a
nonzero element is injective and therefore surjective.  This derives that
\(\kappa(x)\) is finite, but it does not compute ideals, order sites, or build
a database row.

\subsection{Layer 0: Bare Finite Spectrum}

Let \(R\) be a finite commutative ring and \(X=\Spec R\).

\textbf{Local definition.}  The bare \(\Spec R\) layer is the finite
topological/scheme support \(X\), with residue maps \(R\to\kappa(x)\) when
those points and fields have been locally derived or supplied.  It may
discuss supports, opens, specializations, residue fibres, and nilpotent
thickenings in the scheme sense, following the Stacks anchors and AQM-81.  It
has no Pauli operators, Hilbert space, stabilizer Hamiltonian, translation
action, or Haah generator map.  Therefore it is not a Haah code.

This permissive layer is already enough to prevent a common mistake.  A
nilpotent class in \(R\) can be part of the local scheme structure while
having zero image in every residue field.  AQM-34, AQM-35, AQM-36, AQM-37,
AQM-81, and AQM-82 record this distinction for polynomial quotients over
\(\mathbb F_3\): repeated factors change the thickening, not the number of
reduced residue sites.

\subsection{Layer 1: Pre-Haah Module Data}

\textbf{Local definition, from convention \texttt{(ba)}.}  A pre-Haah datum
over \(R\) is an algebraic package \((G,P,\lambda,\sigma)\), where \(G\) and
\(P\) are specified \(R\)-modules, \(\lambda\) is a specified alternating or
symplectic form when chosen, and \(\sigma:G\to P\) is a specified generator
map.  If an adjoint convention is also part of the package, one may write
\(\epsilon=\sigma^\dagger\lambda\).  The quotient
\(\ker(\epsilon)/\operatorname{im}(\sigma)\) is only a logical-defect module,
not automatically a physical logical-Pauli group.

\textbf{Proposal.}  For finite rings, this is the broadest Haah-like grammar:
it keeps the lessons of Haah's matrix formalism, namely generator maps,
commutation tests, and defect quotients, while refusing to add Weyl operators
until a finite nondegenerate commutator pairing and phase character have been
named.  The module package may be interpreted sheaf-theoretically on
\(\Spec R\), but this still does not identify Zariski points with lattice
translation sites.  That separation is exactly the warning in AQM-91 and
convention \texttt{(az)}: coordinate rings and translation group algebras can
look like Laurent rings while playing different roles.

\subsection{Layer 2: Residue Weyl-Haah Data}

\textbf{Local definition.}  Let \(S\subset X\) be a selected finite set of
locally certified points.  Fix an integer \(q\ge1\).  The residue
Weyl-Haah phase module is
\[
  P_{\rm res}(S,q)
  =
  \bigoplus_{x\in S}
  \left(\kappa(x)^q\oplus(\kappa(x)^q)^\vee\right).
\]
For \(e_x=(a_x,\alpha_x)\) and \(f_x=(b_x,\beta_x)\), use
\(\omega_x(e_x,f_x)=\alpha_x(b_x)-\beta_x(a_x)\) and assemble the orthogonal
direct sum over \(S\).  The Weyl carrier is
\[
  \mathcal H_{\rm res}(S,q)
  =\bigotimes_{x\in S}\ell^2(\kappa(x)^q).
\]
This is the finite-field Weyl step of AQM-78 applied sitewise to the
residue-site workflow of AQM-81; it generalizes the finite-set symplectic
label grammar reviewed in AQM-71.

This layer is still not a Haah code.  It becomes a finite residue
Weyl-Haah datum only after a selected generator map
\[
  \sigma_{\rm res}:G_{\rm res}\longrightarrow P_{\rm res}(S,q)
\]
has been supplied and checked to be isotropic with respect to the displayed
pairing, with phases handled by the local Weyl convention.  AQM-22, AQM-40,
and AQM-78 all make the same distinction: the ambient Weyl system is not yet
a stabilizer code.

\subsection{Layer 3: Thickened Frobenius Weyl-Haah Data}

\textbf{Local definition with blocker.}  A thickened finite-ring
Weyl-Haah datum replaces the residue vector spaces by finite ring modules
only when the required self-duality mechanism is already locally sourced:
a finite Frobenius-ring generating character as in AQM-31 and
\path{references/heisenberg_weil/SOURCES.md}, or the explicit
top-coefficient chain-ring character proved in AQM-36 and used in AQM-37.
For a certified finite ring or local factor \(A\), the available model has
\(P_{\rm thick}(A,q)=A^q\oplus A^q\) and
\(\mathcal H_{\rm thick}(A,q)=\ell^2(A^q)\), with translations and phases
defined through the chosen generating character.

For the AQM-36 chain rings
\[
  A=\mathbb F_3[t]/(\pi^e),
\]
the character is the top-coefficient character
\[
  \chi(a)=
  \exp\!\left(2\pi i\,\Tr_{K/\mathbb F_3}(a_{e-1})/3\right),
  \qquad K=\mathbb F_3[t]/(\pi),
\]
and AQM-36 proves it is generating, proves the Weyl product and commutator
laws, and proves the operator-basis statement.  AQM-37 applies this to
\(\mathbb F_3[t]/(t^n-1)\).

The block is part of the definition.  Arbitrary finite commutative rings are
not promoted to this layer in this report.  AQM-78 says that a general
finite-ring Weyl step needs the pairing and character data first; AQM-82 and
convention \texttt{(an)} record that the current finite-ring helpers emit
\texttt{blocked} rows when a certified generating character or recognized
policy is absent.  In particular, this shard does not claim that every finite
commutative ring is Frobenius and does not claim a nondegenerate
\(R^q\oplus R^q\) Weyl label system for arbitrary \(R\).

\subsection{Layer 4: Translation Haah Data}

\textbf{Local definition, from convention \texttt{(ba)}.}  A translation
Haah datum adds a group algebra, an involution or antipode, and the
identity-coefficient trace convention used by Haah's Laurent-polynomial
formalism.  This is the layer of finite Laurent matrices, local translation
constraints, and excitation maps \(\epsilon=\sigma^\dagger\lambda\).
AQM-91 constructs the current finite periodic bridge: choose finite torus
steps, form a quotient of the translation group algebra, put Weyl labels on
the chosen edge module, and compare the AQM-90 square complex with Haah's
toric generator matrix.  The translation basis is group data.  It is not the
same thing as the Zariski point set of \(\Spec R\) unless a separate finite
support bridge identifies them.

\subsection{Permissive-to-Restrictive Ledger}

\begin{center}
\scriptsize\setlength{\tabcolsep}{3pt}
\begin{tabular}{p{0.18\linewidth}p{0.25\linewidth}p{0.22\linewidth}p{0.27\linewidth}}
\toprule
Layer & Added input & Output allowed & Blocker preserved \\
\midrule
Bare \(\Spec R\) &
Finite commutative ring and locally supplied spectrum/residue data. &
Supports, opens, fibres, nilpotent thickenings. &
No Pauli labels, no Hilbert space, no code. \\
Pre-Haah module &
Specified \(R\)-modules \(G,P\), optional \(\lambda\), and
\(\sigma:G\to P\). &
Formal \(\epsilon=\sigma^\dagger\lambda\) and defect module. &
No physical logical group without Weyl data. \\
Residue Weyl-Haah &
Selected finite sites \(S\), residue fields \(\kappa(x)\), and
finite-field duals. &
\(P_{\rm res}(S,q)\), \(\mathcal H_{\rm res}(S,q)\), and possible
isotropic residue generators. &
Nilpotents are invisible except through their effect on residue sites. \\
Thickened Frobenius Weyl &
Certified generating character or sourced Frobenius-ring mechanism. &
\(A^q\oplus A^q\) Weyl labels and \(\ell^2(A^q)\) for certified \(A\). &
Arbitrary finite rings remain \texttt{blocked}. \\
Translation Haah &
Group algebra, antipode/involution, identity-coefficient trace, local
translation generators. &
Laurent matrices, finite-difference excitation maps, periodic logical
quotients. &
Not supplied by \(\Spec R\) alone. \\
\bottomrule
\end{tabular}
\end{center}

\subsection{Example Ledger}

\begin{center}
\scriptsize\setlength{\tabcolsep}{3pt}
\begin{tabular}{p{0.18\linewidth}p{0.24\linewidth}p{0.25\linewidth}p{0.25\linewidth}}
\toprule
Ring or role & Residue layer & Thickened/Frobenius layer & Haah-style lesson \\
\midrule
\(\mathbb F_q\) &
One spectrum point and one residue field site; AQM-22, AQM-78, and AQM-81
give the finite-field Weyl carrier \(\ell^2(\mathbb F_q^q)\) after \(q\) is
chosen. &
No nilpotent thickening in this row.  Finite-field Pauli data are sourced,
but stabilizers are extra choices. &
Bare field geometry gives a site, not a generator matrix. \\
\(k^n\) with \(k=\mathbb F_q\) &
AQM-40 and AQM-83 give \(n\) discrete residue sites with carrier
\(\ell^2(k)^{\otimes n}\) for \(q=1\). &
No thickened layer is imported for the reduced product-field row. &
Tensor factors come from product-spectrum residues; Haah translation still
needs group action and incidence data. \\
\(\mathbb Z/6\mathbb Z\) or mixed residues &
AQM-23 and AQM-81 give residues \(\mathbb F_2\) and \(\mathbb F_3\), hence
\(\mathbb F_2^2\oplus\mathbb F_3^2\) for \(q=1\). &
AQM-80/AQM-82 keep the thickened-Frobenius row blocked without a certified
generating character policy for this mixed row. &
Mixed residue labels are a direct sum of finite abelian phase groups, not one
common-field Laurent module. \\
\(\mathbb F_3[\epsilon]/(\epsilon^2)\) &
AQM-81 records one reduced \(\mathbb F_3\) residue site. &
AQM-36/AQM-82 certify the top-coefficient character and the carrier
\(\ell^2(\mathbb F_3[\epsilon]/(\epsilon^2))\). &
The nilpotent is an internal jet coordinate at the point, not an extra
translation site. \\
\(\mathbb F_3[t]/(t^n-1)\) &
AQM-35 gives the reduced sites from the squarefree part \(t^m-1\), with
reduced dimension \(3^m\). &
AQM-37 gives the certified thickened dimension \(3^n\) for the local
chain-ring factors. &
As a group-algebra quotient it may be used as translation data only after
that role, antipode, trace, and generators are named. \\
Arbitrary finite ring \(R\) with no certified character &
Bare \(\Spec R\) and selected residue sites may be discussed once locally
derived. &
\texttt{blocked}: no nondegenerate \(R^q\oplus R^q\) Weyl layer is asserted. &
The pre-Haah module grammar is allowed, but physical Weyl or code claims wait
for pairing, character, and generator evidence. \\
\bottomrule
\end{tabular}
\end{center}

\subsection{Conclusion}

Finite rings teach a separation principle.  \(\Spec R\) supplies supports,
residue fibres, opens, and nilpotent thickenings.  Haah-like modules require
additional algebraic control data: a generator map, a commutator form, an
adjoint or excitation-map convention, and usually an isotropy check.  Weyl
operators require still more: finite-field residue duals, or a certified
Frobenius/generating-character mechanism for thickened ring labels.
Translation-invariant Haah codes require a further group-algebra role,
antipode or involution, identity-coefficient trace, and local incidence or
translation-generator data.

Thus the safe finite-ring bridge is layered rather than automatic:
\[
  \Spec R
  \leadsto
  \text{pre-Haah module data}
  \leadsto
  \text{residue or certified thickened Weyl data}
  \leadsto
  \text{translation Haah data only after a named bridge}.
\]
Any missing arrow is a blocker, not a convention silently filled in by the
word ``finite''.
